农业机器人通常依赖作物行检测、可通行区域分割与规则后处理构成导航感知链路，二维预测误差易在几何拟合和坐标转换中累积，且不同作物与平台需要重复设计感知接口。本文提出 AgriOcc，一种由正前、左前和右前三路 RGB 图像直接生成机器人前方三维语义占用的端到端框架。AgriOcc 在机器人坐标系下统一预测自由空间、作物、土壤地面、可通行区域、其他植被和其他障碍物，为通行性分析、避障和导航提供统一的三维环境表征。针对农业场景中跨视角可见性不均、三维体素边界和植株结构易被统一解码器平滑、以及远场计算冗余三个关键问题，本文提出几何可见锚点可变形注意力（GVAD）、农业结构自适应混合专家三维占用解码器（Agri-AMoE）和近远场 BEV 查询布局。GVAD 将可靠可见区域压缩为几何锚点，并与局部可变形采样协同建模；Agri-AMoE 将最终 BEV 表征提升为三维体积，以通道/空间显著性和梯度能量路由选择局部、垂直及大范围上下文专家；近远场布局将计算资源优先分配给近场精细结构。基于 UE5 构建的 AgriOcc-Dataset，完整模型在三相机、$512\times288$ 输入及 $100\times100\times25$ 体素网格设置下达到 56.99\% mIoU，较单帧 TSA 基线提升 1.14 个百分点。将 FarmSim 预训练表征适配至 ORAD-3D 的异构类别空间后，模型在真实越野占用评测中保持有效迁移。AgriOcc 为农业移动机器人构建了兼具几何一致性、语义完整性与计算效率的三维环境接口。
